\scanpage{127}
\mspage{63}
\margin{À nouveau se pose la question si $N_V=\operatorname{Norm}(I_V)$. On voudrait que les germes de sections de $\mathcal E_L^{\bar L}$ sur \deleted{$\mathcal E_K$} \inserted{$\mathcal E_{K_1}^{\bar K}$} soient associés \unsure{aux} pts du normalisé $\unsure{\tilde{\mathfrak X}|\tilde{\mathcal Z}}_{L/K}$ de $\unsure{\mathfrak X|\mathcal Z}_{L/K}$ dans $\bar L$ --- qu'elles forment un fibré sur $\unsure{\tilde{\mathfrak X}|\tilde{\mathcal Z}}_{L/K}$, à fibres les $H^1(\deleted{\mathcal E_K}\inserted{\mathcal E_{K_1}}^{\unsure{G|\bar K}},I_V)\ldots$ Ceci est associé à la formule (pour une telle \uline{section} \unsure{$s$} de $\mathcal E_L$ sur $\mathcal E_K$) $\mathcal E_L^{(\mathcal E_K^0)}\simeq I_V$ --- lié à la conjecture que pour $\forall$ \unsure{s-groupe} ouvert de $N_V^0$, ou \unsure{simplt} $\forall$ \unsure{s-groupe} de $N_V$ dont l'image dans $N_V/I_V$ soit \unsure{ouvert}, on a \deleted{\unsure{Centr|tout}} $(\mathcal E_L^{\bar L})^{\unsure{\bar\Gamma|\Gamma}}=I_V$ (ceci pour $\forall$ anneau de valuation $V$ de $L/K$, \unsure{présumé} de corps rés. alg./$K$). En fait, la question est purement géométrique !\ldots}
\unsure{en|à} cette condition --- en fait la section provient d'une section du groupe de décomposition $N_V$ au-dessus de $\mathcal E_K$.

\noindent b) Soit $U$ un modèle ``élémentaire'' de $L$ sur $K$, anabélien. Alors tte section de $\mathcal E_U^{\bar L}$ sur $\mathcal E_K^{\bar K}$ se relève en une section de $\mathcal E_L^{\bar L}$ sur $\mathcal E_K^{\bar K}$.

Elle peut donc se décrire en termes d'un anneau de valuation $V$ de $L/K$, et qd $V$ a un centre sur $U$, \unsure{on n'a qu'à} se \unsure{concentrer} au cas d'un pt $x\in U$ rationnel sur $K$, et de la section définie par $x$. Dans ce cas (\unsure{prenant|posant} tjrs $\Gamma=\mathcal E_K^{\bar K}$) on conjecture encore, bien sûr
\[
(32)\qquad (\mathcal E_L^{\bar L})^{\Gamma^0}=(1)
\]
(et on peut y remplacer $\Gamma^0$ par $\Gamma^0\unsure{\ne1)|\ne)|\subset}\deleted{\unsure{\Gamma}}$
--- mais en fait ceci (en vertu de $\operatorname{Centr}_\Gamma(\Gamma^0)=(1)$) équivaut à la relation analogue \deleted{$\pi_{L\unsure{,|/}K}^{\bar L}$}
\[
(33)\qquad (\pi_{L\unsure{,|/}K}^{\bar L})^{\Gamma^0}=(1)
\]
qui \unsure{se} déduit d'ailleurs, par dévissage,
\margin{126}
